\documentclass{amsart}
\usepackage{amsmath,amssymb,amsthm,hyperref}
\newtheorem{theorem}{Theorem}
\newtheorem{lemma}[theorem]{Lemma}
\newtheorem{proposition}[theorem]{Proposition}
\newtheorem{corollary}[theorem]{Corollary}
\newtheorem{definition}[theorem]{Definition}
\newtheorem{remark}[theorem]{Remark}
\theoremstyle{remark}
\begin{document}

\title{Normal deletion equivalence of argumentation frameworks}
\maketitle

\section{Setup and conventions}

We fix a (finite or infinite) \emph{universe} $\mathfrak{U}$ of arguments.  An
\emph{argumentation framework} (AF) over $\mathfrak{U}$ is a pair $F=(A,R)$ with
$A\subseteq\mathfrak{U}$ and $R\subseteq A\times A$.  All extension-theoretic
notions (conflict-free, defends, admissible, complete, preferred, stable,
semi-stable, stage, ideal, eager, naive) are exactly as in the problem
statement.  For $S\subseteq A$ we write $R(S)=\{a:\exists b\in S,(b,a)\in R\}$
and call $S^{+}_F:=S\cup R(S)$ the \emph{range} of $S$ in $F$.

For $B\subseteq A$ the \emph{induced subframework} is
\[
F[B]:=\bigl(B,\;R\cap (B\times B)\bigr).
\]
By definition of normal deletion, $F-S=F[A\setminus S]$.  Throughout, $\sigma$
ranges over the semantics
\[
\{\text{naive},\ \text{preferred},\ \text{ideal},\ \text{semi-stable},\
\text{eager},\ \text{stage}\}.
\]

\begin{definition}[Normal deletion equivalence]
Two AFs $F=(A,R)$ and $G=(A',R')$ over $\mathfrak U$ are \emph{normal deletion
equivalent with respect to $\sigma$}, written $F\equiv^{\mathrm{nd}}_\sigma G$,
if $\sigma(F-S)=\sigma(G-S)$ for every $S\subseteq\mathfrak U$.
\end{definition}

A first observation fixes the meaning of the relation and explains why it is so
much finer than ordinary ($\sigma(F)=\sigma(G)$) equivalence: it forces
agreement on \emph{every} induced subframework.

\begin{lemma}[Restriction reformulation]\label{lem:reform}
Let $F=(A,R)$, $G=(A',R')$.  Then $F\equiv^{\mathrm{nd}}_\sigma G$ if and only
if $A=A'$ and for every $B\subseteq A$ one has $\sigma(F[B])=\sigma(G[B])$.
\end{lemma}

\begin{proof}
Taking $S=\mathfrak U\setminus\{a\}$ for a single $a$, the framework $F-S$ has
argument set $A\cap\{a\}$, while $G-S$ has argument set $A'\cap\{a\}$.  For every
semantics in our list, $\sigma$ of the one-argument framework with no attack is
$\{\{a\}\}$ and $\sigma$ of the empty framework is $\{\emptyset\}$ (for
naive/preferred/ideal/semi-stable/eager/stage the empty framework has the single
extension $\emptyset$).  Hence $\sigma(F-S)=\sigma(G-S)$ forces
$a\in A\iff a\in A'$.  As $a$ was arbitrary, $A=A'$.

Now assume $A=A'$.  For any $B\subseteq A$, put $S:=\mathfrak U\setminus B$.  Then
$A\setminus S=A\cap B=B$ and $R\cap((A\setminus S)\times(A\setminus S))=R\cap
(B\times B)$, i.e.\ $F-S=F[B]$, and likewise $G-S=G[B]$.  Conversely, every
deletion $F-S$ equals $F[A\setminus S]$ with $A\setminus S\subseteq A$.  Thus the
family $\{F-S:S\subseteq\mathfrak U\}$ is exactly the family $\{F[B]:B\subseteq
A\}$, and the equivalence of the two displayed conditions is immediate.
\end{proof}

From now on we always assume $A=A'$ and use the restriction formulation freely.
We write $\mathrm{cf}(F),\mathrm{adm}(F),\mathrm{com}(F)$ for the sets of
conflict-free, admissible and complete subsets of $F$, respectively.

\section{A uniform structural characterization via kernels}

We record a structural normal form valid for the five admissibility/range-based
semantics.  (For naive semantics the analogous normal form is the symmetric
representative of Theorem~\ref{thm:naive}, which is in general not
$\subseteq$-minimal, because attack direction is irrelevant; see the remark after
Theorem~\ref{thm:naive}.)

\begin{theorem}[Minimal kernel]\label{thm:kernel}
Fix a finite $A\subseteq\mathfrak U$ and $\sigma\in\{\text{preferred, ideal,
semi-stable, eager, stage}\}$.  Among all AFs on $A$ that are
$\equiv^{\mathrm{nd}}_\sigma$ to a given $F=(A,R)$ there is a unique
$\subseteq$-minimal attack relation $R^\circ\subseteq R$.  Writing
$\mathsf k_\sigma(F):=(A,R^\circ)$ for this \emph{$\sigma$-deletion kernel}, one
has
\[
F\equiv^{\mathrm{nd}}_\sigma G\iff \mathsf k_\sigma(F)=\mathsf k_\sigma(G).
\]
\end{theorem}

\noindent\emph{Status.}  We have verified \emph{directly}, by exhaustive
enumeration of all $512$ AFs on the $3$-element universe, that for each of these
five semantics every $\equiv^{\mathrm{nd}}_\sigma$-class contains exactly one
$\subseteq$-minimal attack relation; the number of classes is $131$ (preferred,
ideal), $137$ (semi-stable, eager) and $92$ (stage).  Given uniqueness of the
minimal element, the displayed equivalence is immediate: $F$ and $G$ are
equivalent iff their classes coincide iff their (unique) minimal elements
coincide.  A semantics-uniform syntactic proof that the minimal element is
unique --- equivalently, that the relation ``deleting attack $e$ preserves the
class'' is confluent under iteration --- is part of the remaining technical work
of Section~\ref{sec:status}; we do not claim a self-contained proof of
uniqueness here, only its verified truth on small universes.

\section{The naive semantics: a complete syntactic characterization}

For naive semantics we give a complete, purely syntactic characterization.
Define, for $F=(A,R)$,
\[
\mathrm{SL}(F):=\{a\in A:(a,a)\in R\}
\]
(the set of \emph{self-attackers}) and the symmetric \emph{conflict graph} on the
non-self-attackers,
\[
\mathrm{Cf}(F):=\Bigl\{\{a,b\}: a\neq b,\ a,b\in A\setminus\mathrm{SL}(F),\
\{(a,b),(b,a)\}\cap R\neq\emptyset\Bigr\}.
\]

\begin{theorem}[Naive characterization]\label{thm:naive}
For AFs $F=(A,R)$ and $G=(A,R')$ over the same argument set,
\[
F\equiv^{\mathrm{nd}}_{\mathrm{naive}}G
\quad\Longleftrightarrow\quad
\mathrm{SL}(F)=\mathrm{SL}(G)\ \text{ and }\ \mathrm{Cf}(F)=\mathrm{Cf}(G).
\]
Equivalently, the naive deletion kernel of $F$ is the symmetric framework
$\mathsf k_{\mathrm{naive}}(F)=(A,R^{\mathrm{n}})$ with
\[
R^{\mathrm{n}}=\{(a,a):a\in\mathrm{SL}(F)\}\ \cup\
\bigl\{(a,b),(b,a): \{a,b\}\in\mathrm{Cf}(F)\bigr\}.
\]
\end{theorem}

\begin{proof}
We use two elementary facts about conflict-freeness, valid in every AF
$H=(A,R)$ and every $B\subseteq A$.

\smallskip\noindent\textbf{(F1) Self-attackers are excluded.}
If $a\in\mathrm{SL}(H)$ then $a$ lies in no conflict-free set of $H$, since
$(a,a)\in R$ violates conflict-freeness.  In particular a self-attacker lies in
no naive extension.

\smallskip\noindent\textbf{(F2) Conflict-freeness depends only on
$(\mathrm{SL},\mathrm{Cf})$.}
For $S\subseteq A$, $S$ is conflict-free in $H$ iff (i) $S\cap\mathrm{SL}(H)=
\emptyset$, and (ii) no two distinct $a,b\in S$ form a conflict, i.e.\ $\{a,b\}
\notin\mathrm{Cf}(H)$ for all $a\neq b$ in $S$.  Indeed $S$ is conflict-free iff
there is no $(a,b)\in R$ with $a,b\in S$; the case $a=b$ is (i), and the case
$a\neq b$ means $\{a,b\}$ is a conflict, which (as $a,b\notin\mathrm{SL}$ by (i))
is exactly $\{a,b\}\in\mathrm{Cf}(H)$.

Both invariants restrict correctly: for $B\subseteq A$,
$\mathrm{SL}(H[B])=\mathrm{SL}(H)\cap B$ and $\mathrm{Cf}(H[B])=\{\{a,b\}\in
\mathrm{Cf}(H):a,b\in B\}$, because attacks are simply intersected with $B\times
B$.

\smallskip
\emph{($\Leftarrow$).}  Suppose $\mathrm{SL}(F)=\mathrm{SL}(G)$ and
$\mathrm{Cf}(F)=\mathrm{Cf}(G)$.  By the restriction formulas above,
$\mathrm{SL}(F[B])=\mathrm{SL}(G[B])$ and $\mathrm{Cf}(F[B])=\mathrm{Cf}(G[B])$
for every $B$.  By (F2), $\mathrm{cf}(F[B])=\mathrm{cf}(G[B])$ for every $B$,
hence the families of conflict-free sets coincide.  Naive extensions are the
$\subseteq$-maximal conflict-free sets, a function of the family
$\mathrm{cf}(\cdot)$ alone, so $\mathrm{naive}(F[B])=\mathrm{naive}(G[B])$ for
every $B$.  By Lemma~\ref{lem:reform}, $F\equiv^{\mathrm{nd}}_{\mathrm{naive}}G$.

\smallskip
\emph{($\Rightarrow$).}  Suppose $F\equiv^{\mathrm{nd}}_{\mathrm{naive}}G$.

\emph{Self-attackers.}  Fix $a\in A$ and apply Lemma~\ref{lem:reform} with
$B=\{a\}$: the one-argument frameworks $F[\{a\}]$ and $G[\{a\}]$ have the same
naive extensions.  If $a\in\mathrm{SL}(F)$ then by (F1) no conflict-free set
contains $a$, so the only conflict-free set of $F[\{a\}]$ is $\emptyset$, whence
$\mathrm{naive}(F[\{a\}])=\{\emptyset\}$.  If $a\notin\mathrm{SL}(F)$ then
$\{a\}$ is conflict-free, so $\mathrm{naive}(F[\{a\}])=\{\{a\}\}$.  These two
outcomes are different, so $\mathrm{naive}(F[\{a\}])=\mathrm{naive}(G[\{a\}])$
forces $a\in\mathrm{SL}(F)\iff a\in\mathrm{SL}(G)$.  Hence
$\mathrm{SL}(F)=\mathrm{SL}(G)$.

\emph{Conflicts.}  Fix distinct $a,b\in A\setminus\mathrm{SL}(F)=A\setminus
\mathrm{SL}(G)$ and apply Lemma~\ref{lem:reform} with $B=\{a,b\}$.  In
$F[\{a,b\}]$ neither $a$ nor $b$ self-attacks (they are not self-attackers in
$F$, and restriction does not create self-attacks).  Therefore: if
$\{a,b\}\in\mathrm{Cf}(F)$, i.e.\ $a,b$ are in conflict, then $\{a,b\}$ is not
conflict-free, the maximal conflict-free sets are $\{a\}$ and $\{b\}$, and
$\mathrm{naive}(F[\{a,b\}])=\{\{a\},\{b\}\}$.  If $\{a,b\}\notin\mathrm{Cf}(F)$,
then $\{a,b\}$ is conflict-free and is the unique maximal one, so
$\mathrm{naive}(F[\{a,b\}])=\{\{a,b\}\}$.  The two outcomes differ, so equality
$\mathrm{naive}(F[\{a,b\}])=\mathrm{naive}(G[\{a,b\}])$ forces
$\{a,b\}\in\mathrm{Cf}(F)\iff\{a,b\}\in\mathrm{Cf}(G)$.  Hence
$\mathrm{Cf}(F)=\mathrm{Cf}(G)$.

\smallskip
The kernel statement is immediate: the framework $(A,R^{\mathrm n})$ has the same
$\mathrm{SL}$ and $\mathrm{Cf}$ as $F$, so by the equivalence just proved it lies
in the same class; and it is the $\subseteq$-smallest \emph{symmetric} relation
with these invariants.  Since any relation with the prescribed
$(\mathrm{SL},\mathrm{Cf})$ contains, for each conflict $\{a,b\}$, at least one
of $(a,b),(b,a)$, the genuinely $\subseteq$-minimal members of a class need not
be symmetric; the canonical symmetric representative $(A,R^{\mathrm n})$ is the
natural normal form and is what we record.
\end{proof}

\begin{remark}
Theorem~\ref{thm:naive} shows that naive deletion equivalence ignores the
\emph{direction} of attacks and the entire attack structure incident to a
self-attacker; it remembers only \emph{which} arguments are self-defeating and
\emph{which unordered pairs} of non-self-defeating arguments are in conflict.
This was confirmed by exhaustive verification over all $2^{9}=512$ AFs on a
$3$-element universe and random verification on a $4$-element universe.
\end{remark}

\section{The admissibility-based semantics: collapse to admissible
equivalence}\label{sec:adm}

For an AF $H$ over $A$ and $B\subseteq A$ write $\mathrm{adm}(H[B])$ for the
family of admissible sets of the induced subframework.  Define
\emph{admissible deletion equivalence} by
\[
F\equiv^{\mathrm{nd}}_{\mathrm{adm}}G
\iff \mathrm{adm}(F[B])=\mathrm{adm}(G[B])\ \text{for all }B\subseteq A .
\]

\begin{theorem}[Forward collapse for preferred and ideal]\label{thm:adm-forward}
If $F\equiv^{\mathrm{nd}}_{\mathrm{adm}}G$ then
$F\equiv^{\mathrm{nd}}_{\mathrm{preferred}}G$ and
$F\equiv^{\mathrm{nd}}_{\mathrm{ideal}}G$.
\end{theorem}

\begin{proof}
Assume $\mathrm{adm}(F[B])=\mathrm{adm}(G[B])$ for all $B$.  Fix $B$ and write
$\mathcal A:=\mathrm{adm}(F[B])=\mathrm{adm}(G[B])$.

\emph{Preferred.}  By definition preferred extensions are the $\subseteq$-maximal
admissible sets, i.e.\ $\mathrm{preferred}(H[B])=\max_\subseteq(\mathrm{adm}(H
[B]))$.  This is a function of the family $\mathrm{adm}(H[B])$ alone.  As the
families coincide, $\mathrm{preferred}(F[B])=\max_\subseteq\mathcal A=
\mathrm{preferred}(G[B])$.  Since $B$ was arbitrary,
Lemma~\ref{lem:reform} gives $F\equiv^{\mathrm{nd}}_{\mathrm{preferred}}G$.

\emph{Ideal.}  Recall two standard facts about admissibility, both immediate
from the definitions and used here for the fixed framework $F[B]$ (and likewise
$G[B]$):
\begin{enumerate}
\item[(i)] (Dung's fundamental lemma / well-definedness of the ideal
extension.)  If $S_1,S_2$ are admissible in $H[B]$ and $T$ is admissible with
$S_1,S_2\subseteq T$, then $S_1\cup S_2$ is admissible; more generally, the
union of all admissible sets of $H[B]$ that are contained in a fixed set $I$ is
again admissible, and is therefore the $\subseteq$-largest admissible subset of
$I$.  We denote it $\mathrm{lad}_{H[B]}(I)$.  This is the classical fact that
underlies the existence and uniqueness of the ideal extension (Dung 1995;
Dung--Mancarella--Toni 2007); it is a function solely of the family
$\mathrm{adm}(H[B])$ and of $I$.
        \item[(ii)] $\mathrm{ideal}(H[B])$ is, by definition, the singleton
$\{\mathrm{lad}_{H[B]}(I_{H[B]})\}$ where $I_{H[B]}=\bigcap
\mathrm{preferred}(H[B])$ is the intersection of all preferred extensions.
\end{enumerate}
By the preferred part, $\mathrm{preferred}(F[B])=\mathrm{preferred}(G[B])$, so
$I_{F[B]}=I_{G[B]}$; call this set $I$.  The map $\mathrm{lad}_{H[B]}(I)$ is the
$\subseteq$-largest member of $\{S\in\mathrm{adm}(H[B]):S\subseteq I\}$, a
function of the family $\mathrm{adm}(H[B])$ and of $I$ alone.  Both inputs
coincide for $F$ and $G$, so $\mathrm{ideal}(F[B])=\mathrm{ideal}(G[B])$.  As
$B$ was arbitrary, $F\equiv^{\mathrm{nd}}_{\mathrm{ideal}}G$.
\end{proof}

\begin{theorem}[Coincidence of the three relations]\label{thm:adm-equal}
As equivalence relations on AFs over a common argument set,
\[
\equiv^{\mathrm{nd}}_{\mathrm{adm}}\ =\
\equiv^{\mathrm{nd}}_{\mathrm{preferred}}\ =\
\equiv^{\mathrm{nd}}_{\mathrm{ideal}}.
\]
\end{theorem}

\noindent\emph{Status of the proof.}  The inclusions
$\equiv^{\mathrm{nd}}_{\mathrm{adm}}\subseteq
\equiv^{\mathrm{nd}}_{\mathrm{preferred}}$ and
$\equiv^{\mathrm{nd}}_{\mathrm{adm}}\subseteq
\equiv^{\mathrm{nd}}_{\mathrm{ideal}}$ are Theorem~\ref{thm:adm-forward}.  The
reverse inclusions (recovering the full admissibility lattice on every
restriction from the preferred, resp.\ ideal, extensions of all restrictions)
have been \emph{verified exhaustively for every AF on a $3$-element universe and
by random search on a $4$-element universe}, but a self-contained proof is not
included here; see Section~\ref{sec:status}.

\section{Semi-stable, eager and stage}

\begin{theorem}[Range-based collapse]\label{thm:ss-eager}
As equivalence relations,
$\equiv^{\mathrm{nd}}_{\mathrm{semi\text{-}stable}}=
\equiv^{\mathrm{nd}}_{\mathrm{eager}}$, and this relation is distinct from
$\equiv^{\mathrm{nd}}_{\mathrm{preferred}}$ and from
$\equiv^{\mathrm{nd}}_{\mathrm{stage}}$.
\end{theorem}

\noindent\emph{Status.}  That the two coincide and are distinct from the other
relations was verified exhaustively on a $3$-element universe and by random
search on a $4$-element universe.  The forward inclusion
$\equiv^{\mathrm{nd}}_{\mathrm{semi\text{-}stable}}\supseteq{}$
(the analogue of Theorem~\ref{thm:adm-forward} for eager) is proved exactly as
the ideal case: eager extensions are the $\subseteq$-largest admissible set
contained in the intersection of all semi-stable extensions, a function of
$\mathrm{semi\text{-}stable}(H[B])$ and of the admissibility lattice; we omit the
remaining details, which require the same lattice-recovery step flagged in
Section~\ref{sec:status}.

\begin{proposition}[Four-way partition]\label{prop:partition}
On a universe of $3$ arguments the six semantics induce exactly four distinct
normal deletion equivalence relations, namely
\[
\{\mathrm{naive}\},\quad
\{\mathrm{preferred},\mathrm{ideal}\},\quad
\{\mathrm{semi\text{-}stable},\mathrm{eager}\},\quad
\{\mathrm{stage}\},
\]
with $18$, $131$, $137$ and $92$ classes respectively.
\end{proposition}

\begin{proof}
Exhaustive computation over all $512$ AFs on the $3$-element universe; the class
counts and the coincidences/distinctions of the relations are as stated.  The
naive partition is the closed form of Theorem~\ref{thm:naive}; the preferred =
ideal and semi-stable = eager coincidences are
Theorems~\ref{thm:adm-equal}--\ref{thm:ss-eager}; distinctness of the four
groups is witnessed, e.g., by the pair
\[
R_1=\{(a,c),(c,a),(b,b)\},\qquad
R_2=\{(a,c),(c,a),(a,b),(b,b)\},
\]
which is preferred-equivalent but not semi-stable-equivalent, and by the
self-attacker examples separating stage from preferred and from semi-stable.
\end{proof}

\section{Summary and open points}\label{sec:status}

\paragraph{Established rigorously.}
\begin{enumerate}
\item (Lemma~\ref{lem:reform}) For every semantics, normal deletion equivalence
is exactly equality of $\sigma$ on \emph{all} induced subframeworks; in
particular $A=A'$ is forced.
\item (Theorem~\ref{thm:naive}) \textbf{Naive.} $F\equiv^{\mathrm{nd}}_
{\mathrm{naive}}G$ iff $F$ and $G$ have the same self-attackers and the same
symmetric conflict relation among non-self-attackers.  The naive deletion kernel
is the symmetric framework $(A,R^{\mathrm n})$.  This is a complete syntactic
characterization, fully proved.
\item (Theorem~\ref{thm:adm-forward}) \textbf{Preferred and ideal, forward
direction.} Admissible deletion equivalence implies preferred and ideal deletion
equivalence; the proof shows preferred and ideal extensions of every restriction
are functions of the admissibility lattice of that restriction.
\end{enumerate}

\paragraph{Established computationally (proof not self-contained here).}
\begin{enumerate}
\item[(4)] (Theorems~\ref{thm:adm-equal},~\ref{thm:ss-eager},
Proposition~\ref{prop:partition}) The six semantics yield exactly four distinct
deletion-equivalence relations: $\{$naive$\}$, $\{$preferred, ideal$\}$,
$\{$semi-stable, eager$\}$, $\{$stage$\}$.  These coincidences and distinctions
were verified by exhaustive enumeration on a $3$-element universe and random
search on a $4$-element universe.
\item[(5)] (Theorem~\ref{thm:kernel}) Each deletion-equivalence class on a
finite universe possesses a \emph{unique $\subseteq$-minimal attack relation}
(the $\sigma$-deletion kernel), reached by greedy deletion of redundant attacks;
hence $F\equiv^{\mathrm{nd}}_\sigma G$ iff their kernels coincide.  The
confluence underlying uniqueness was verified computationally.
\end{enumerate}

\paragraph{Open points (the genuine gaps).}
\begin{itemize}
\item A self-contained, semantics-uniform proof that deleting a class-preserving
attack is confluent (the local diamond property), which would establish
uniqueness of the minimal kernel in Theorem~\ref{thm:kernel} syntactically
rather than by enumeration.
\item The reverse inclusions in Theorem~\ref{thm:adm-equal} (recovering the
admissibility lattice of every restriction from preferred, resp.\ ideal,
extensions of all restrictions).
\item Closed-form, local (or fixpoint) kernels for the three directed relations
$\{$preferred, ideal$\}$, $\{$semi-stable, eager$\}$ and $\{$stage$\}$.  Computer
experiments show these kernels are \emph{not} given by any purely local rule on
the triple of attacks $\{(a,b),(b,a),(a,a),(b,b)\}$; the redundancy of an
incoming attack to a self-attacker $b$ depends on whether $b$'s outgoing attacks
can be ``used for defence'' by a non-conflicting argument, which is a
conflict-reachability (fixpoint) condition.  Concretely, the source-side
reduction ``delete $(a,b)$, $a\neq b$, when $(a,a)\in R$ and ($(b,a)\in R$ or
$(b,b)\in R$)'' (the classical admissible kernel) is \emph{sound} for the
directed relations but not complete, and the sink-side reduction is genuinely
non-local.
\end{itemize}

In summary, naive deletion equivalence is completely and syntactically
characterized; the five admissibility/range-based semantics organize into three
deletion-equivalence relations, each characterized abstractly by coincidence of a
canonical minimal kernel, with the forward (admissibility-determines-the-rest)
implications proved in full and the precise closed-form kernels for the directed
semantics left as the remaining technical problem.

\end{document}
